\documentclass[10pt]{article}
\usepackage[margin=20mm]{geometry}
\usepackage{amsmath,amssymb,mathtools,fontspec,unicode-math,xurl,hyperref}
\IfFontExistsTF{DejaVu Serif}{\setmainfont{DejaVu Serif}}{\setmainfont{Latin Modern Roman}}
\setmathfont{latinmodern-math.otf}
\setlength{\parindent}{0pt}\setlength{\parskip}{6pt}\sloppy
\title{Technical review of
Motivic monodromy and \(p\)-adic cohomology theories}
\author{AI technical verification}
\date{}
\begin{document}
\maketitle
Abstract: The reviewed artifact is the Warwick accepted manuscript by Federico Binda, Martin Gallauer and Alberto Vezzani. Its main text was read, with focused checks of the categorical equivalences, weight criterion, comparison argument and Clemens–Schmid construction. One confirmed issue affects a premise explicitly used in the central comparison proof: the stated integrality condition on pure Frobenius modules excludes positive Tate twists allowed by the proposition. This is assessed Major at the level of that derivation, without asserting that the comparison theorem is false. Two further, local issues concern an infinite-product coalgebra and the Hodge exposition. Explicit counterexamples and bounded repairs are given. External foundational theorems and all higher coherence data were not independently reconstructed. Scientific adjudication remains pending.\par
\section*{Identity, source and coverage}
The assigned identifier is M25-0261, in the frozen 2025 batch. The original title is Motivic monodromy and \(p\)-adic cohomology theories; the authors are Federico Binda, Martin Gallauer and Alberto Vezzani. The journal identifier is DOI `10.4171/jems/1634`. The actual 51-page PDF is the Warwick WRAP record 188439 accepted manuscript (accepted 20 October 2024; deposited 7 November 2024). Its repository cover expressly distinguishes it from the version of record. The associated arXiv v2 date, 31 October 2024, is not a byte-equivalence claim. The publisher’s online date is 4 April 2025. These dates do not change the frozen membership.\par
Here PDF page counts the repository cover: printed page \(r\) is PDF page \(r+1\). All main-text pages 2–48 and the bibliography on pages 48–51 were read. The review checks the distinction between ind-nilpotent and arbitrary operators; square-zero and monodromy conventions; the weight-orthogonality criterion; the finite-residue-field comparison and Kummer normalization; and the chain property, degrees and twists of the Clemens–Schmid construction. The accompanying coverage ledger distinguishes reading from independent proof reconstruction. Selected decisive pages were also inspected as rendered source images.\par
No supplementary mathematical file or author code was identified in the bounded searches recorded with the source manifest. This is an availability limitation, not evidence of absence and not a defect. No author code was executed. The search for corrections did not establish a paper-specific correction to the exact accepted manuscript. A later publisher PDF was located but not compared; no conclusion here is transferred to that version without verification.\par
\section*{Framework and checked dependencies}
Section 2 works in a \(\mathbb Q\)-linear, compactly generated stable symmetric monoidal category \(\mathcal C\), with an object \(t\) satisfying \(\operatorname{Alt}^2(t)=0\). The ind-nilpotent condition is the vanishing of the telescope of \(X\to X\otimes t\to X\otimes t^{\otimes2}\to\cdots\). Theorem 2.1 and Propositions 2.32 and 2.38 use the square-zero algebra \(\mathbf 1\oplus t[-1]\) and the direct sum \(\bigoplus_{n\ge0}t^{\otimes n}\). They must not be conflated with the arbitrary-operator extension in Remark 2.40.\par
The shift in Lemma 2.31 changes the symmetry sign, so that \(\operatorname{Sym}^n(t[-1])=0\) for \(n\ge2\). In the direct-sum right-adjoint calculation, mapping a compact zero-section generator into the cofree object identifies the fiber of the shift with the degree-zero summand. This checks the relevant adjunction without commuting an arbitrary infinite product through tensor product.\par
The application to rigid motives is restricted to the stated good-reduction subcategory unless the source supplies the additional hypotheses giving the whole category. Corollary 3.18 requires both underlying-realization compatibility and Kummer normalization. The geometric and arithmetic Frobenius directions are inverses as explained in Remark 2.52; Lemma 4.30 and Remark 4.41 must be read together. The mixed-characteristic construction uses \(\varphi^\sharp\), not the unsuitable naive lift discussed in Remark 4.47.\par
Example 3.31 needs cohomological dimension at most one, different-weight Hom vanishing and a specified direction of Ext vanishing. It does not assume that an entire mixed category is semisimple. Theorem 4.53 uses finite-field purity for every object of the full Chow heart, and then descends to \(k\subset\overline{\mathbb F}_p\). Corollaries 4.54 and 4.58 retain that residue-field restriction; Corollary 4.58 also fixes the stated uniformizer normalization in characteristic zero.\par
\section*{Confirmed findings}
Problem 1 (Tate twists contradict the purity premise used in the comparison).  Definition 2.42(3), PDF page 18 (printed 17), imposes algebraic integrality on every Frobenius eigenvalue of a pure module. Proposition 4.44(2), PDF page 39 (printed 38), asserts purity for \(R^i\Gamma_{\mathrm{rig}}(M)\) for every \(M\in\mathcal H_k\) when \(k\) is finite. With the stated definition and the full heart, this assertion is false.\par
\subsection*{Hypotheses and exact computation}
Take \(k=\mathbb F_2\), \(K=K_0=\mathbb Q_2\) and \[M=\mathbf 1(1)[2]=M(\operatorname{Spec}\mathbb F_2)\{1\}.\] The residue field is finite and perfect, the field \(K\) is complete with rational-rank-one value group, and the point is smooth and proper. Example 3.29(1), PDF page 29 (printed 28), includes all integer twists \(\{r\}=(r)[2r]\), so \(M\in\mathcal H_k\).\par
Example 2.43(2) specifies \(K_0(r)=(K_0,2^{-r}\varphi)\). Remark 4.41 specifies the realization of \(\mathbf 1(-1)\) as \(K_0(-1)\), with Frobenius \(2\varphi\). Exact monoidality and tensor inversion therefore give \[R\Gamma_{\mathrm{rig}}(\mathbf 1(1)[2])\simeq \mathbb Q_2(1)[2].\] Witt Frobenius on \(\mathbb Q_2\) is the identity. The unique nonzero graded piece has eigenvalue \(\lambda=1/2\). Under cohomological indexing it occurs in degree \(-2\); the integrality contradiction is independent of indexing.\par
If \(1/2\) satisfied a monic integral polynomial of degree \(m\), multiplication of that equation by \(2^m\) would yield \[1+2a_{m-1}+4a_{m-2}+\cdots+2^m a_0=0,\qquad a_j\in\mathbb Z,\] which is impossible modulo \(2\). Thus this nonzero piece is not pure of any weight under the printed integrality definition. Its correct spectral weight is nevertheless \(-2\), since \(|1/2|=2^{-2/2}\).\par
\subsection*{Alternative readings, impact and repair}
This is not an OCR artifact: the integrality word, the integer-twist quantifier and the universal purity claim were checked on source images. Reversing Frobenius conventions cannot help after Remark 4.41 has fixed the arithmetic realization of \(\mathbf 1(-1)\). Restricting to effective motives would change the stated heart and would not justify the application to all of \(\mathcal H_k\).\par
The proof of Theorem 4.53, PDF page 41 (printed 40), explicitly invokes Proposition 4.44 for arbitrary \(M\in\mathcal H_k\) before applying Example 3.31 and Corollary 4.38. Thus a demonstrably false universal input occurs in a core comparison derivation, subsequently used in Corollaries 4.54 and 4.58. The assessment is Major for this derivation. It does not assert a counterexample to the resulting Hyodo–Kato comparison or the preceding motivic equivalence. The classification follows the demonstrated dependency, not the length of the repair.\par
The appropriate local repair is to require nonzero algebraic numbers whose complex absolute values are \(q^{w/2}\), without the extra integrality requirement. Tate twisting then sends an eigenvalue \(\lambda\) to \(q^{-r}\lambda\) and weight \(w\) to \(w-2r\). The orthogonality needed in the proof survives: distinct weights have disjoint spectra for the linear operator \(\varphi^a\), so the corresponding Hom group vanishes. The two terms computing \(\mathop{\mathrm{Ext}}^1\) in Remark 2.48 have equal finite \(\mathbb Q_p\)-dimension; injectivity of their differential therefore gives surjectivity as well. Neither step uses integrality. This verifies the local spectral repair and the relevant Hom/Ext calculation, not a fresh proof of the external Weil theorem or every coherence in the complete comparison argument.\par
Problem 2 (The product object is not the asserted ordinary internal coalgebra).  Remark 2.40, PDF page 17 (printed 16), identifies all lax fixed points with comodules over \(\widehat H(t)=\prod_{n\ge0}t^{\otimes n}\). The internal-coalgebra and ordinary tensor interpretation established in §2.5 does not support that identification.\par
\subsection*{A counterexample in the exact categorical setting}
Let \(\mathcal C=D(\mathbb Q)\) and \(t=\mathbb Q[0]\). This is a presentable, compactly generated, \(\mathbb Q\)-linear stable symmetric monoidal category; \(t\) is compact, \(\operatorname{Alt}^2(t)=0\), and \(-\otimes t\) is the identity, which also preserves products. These satisfy the relevant assumptions, including the extra product-preservation condition.\par
For arbitrary endomorphisms the genuine forgetful right adjoint sends \(X\) to \((\prod_{n\ge0}X,\text{left shift})\); its comonad is \(G(X)=\prod_{n\ge0}X\). For ordinary internal comodules the corresponding functor would be \(\widehat H\otimes X\) with \(\widehat H=\prod_{n\ge0}\mathbb Q\). Put \(X=\bigoplus_{j\ge0}\mathbb Qe_j\). Every element of \[(\prod_{n\ge0}\mathbb Q)\otimes X\longrightarrow\prod_{n\ge0}X\] is a sequence lying in one common finite-dimensional subspace: write the input as a finite sum \(\sum_{\ell=1}^r a^{(\ell)}\otimes x_\ell\). Its coordinates all lie in \(\operatorname{span}(x_1,\ldots,x_r)\). The sequence \((e_n)_{n\ge0}\) does not. It occurs as the orbit of \(e_0\) for the admissible lax fixed point \(S(e_j)=e_{j+1}\). Remark 2.40 expressly permits nonnilpotent operators, so this is in its scope.\par
There is also a direct obstruction to the proposed canonical comultiplication. At \(\mathbb Q\), the comonad sends \(a=(a_n)\) to \((a_{i+j})_{i,j\ge0}\). Choose \(a_n=1/(n+1)\). Its matrix has infinite rank: for every nonzero \(c\in\mathbb R^m\), \[\sum_{i,j=0}^{m-1}\frac{c_i c_j}{i+j+1}
 =\int_0^1\left(\sum_{i=0}^{m-1}c_i x^i\right)^2\,dx>0.\] Every finite leading block is nonsingular, whereas every image of \(\widehat H\otimes\widehat H\) is a finite sum of rank-one matrices. The canonical map therefore cannot factor through that ordinary tensor product. This does not rule out unrelated coalgebra structures on the underlying vector space.\par
\subsection*{Scope and supported correction}
Comonadicity of the forgetful functor does not identify a comonad with tensoring by its value on the unit. Keeping coalgebras for the actual comonad \(G\) is a valid replacement in this example. A completed tensor product or a pro-category would require an explicit different statement and its own proof. No such replacement is silently assumed here. The consulted primary discussions [Retained latex: \textbackslash{}cite[\textbackslash{}S3.5.2]\{FG\}] and [Retained latex: \textbackslash{}cite[Chapter 6, \textbackslash{}S2.7.3--2.7.6]\{GR\}] themselves distinguish product/right-lax constructions from the corresponding comonad; they do not supply the missing internal-tensor identification.\par
This is Minor: §2.5 explicitly says it is not used later, and the main ind-nilpotent theorem uses a direct sum instead of this product. No central dependency on the false arbitrary-operator extension was established. Its peripheral role, rather than ease of deletion, determines the assessment.\par
Problem 3 (Two incorrect descriptions in the Hodge realization paragraph).  Remark 4.62, PDF page 44 (printed 43), calls the category of rational polarizable mixed Hodge structures semisimple and gives an off-diagonal vanishing criterion for intersections of Hodge filtrations. Both assertions fail as written.\par
\subsection*{A nonsplit object with polarizable weight-graded pieces}
Let \(E_\mathbb Q=\mathbb Qe_0\oplus\mathbb Qe_1\). Define \[W_mE_\mathbb Q=\begin{cases}
 0&m<-2,\\ \mathbb Qe_1&-2\le m<0,\\ E_\mathbb Q&m\ge0,
 \end{cases}
 \qquad
 F^pE_\mathbb C=\begin{cases}
 0&p\ge1,\\
 \mathbb C(e_0+i e_1)&p=0,\\
 E_\mathbb C&p\le-1 .
 \end{cases}\] The induced graded pieces are of types \((-1,-1)\) and \((0,0)\). They are rank-one Tate structures and admit polarizations. Thus \(E\) is a graded-polarizable mixed Hodge structure, in the standard meaning of the target category; see the separately consulted [Retained latex: \textbackslash{}cite[Example 1.4(ii), p.\textasciitilde{}9]\{Drew\}], which also records homological dimension one in Remark 1.5(ii).\par
The subobject \(\mathbb Qe_1\) and quotient \(\mathbb Q\bar e_0\) give an exact sequence \[0\longrightarrow\mathbb Q(1)\longrightarrow E\longrightarrow\mathbb Q(0)
 \longrightarrow0.\] A rational section must send \(1\) to \(e_0+a e_1\) for \(a\in\mathbb Q\). Preservation of \(F^0\) forces \(a=i\), which is impossible. Hence the sequence does not split and the mixed category is not semisimple. This objection does not misapply the familiar semisimplicity of pure polarizable Hodge structures of a fixed weight.\par
\subsection*{The adjacent intersection formula}
The same paragraph sets \(V_\mathbb C^{p,q}=F^pV_\mathbb C\cap\overline{F^qV_\mathbb C}\) and requires its vanishing for \(p+q\ne n\) to express weight \(n\). The source image confirms the conjugation and the quantifier. Its own next example \(\mathbb Q(1)\) has \(F^{-1}=V_\mathbb C\), hence \(F^{-2}=V_\mathbb C\); therefore \[F^{-2}V_\mathbb C\cap\overline{F^{-2}V_\mathbb C}=V_\mathbb C\ne0,\qquad
 (-2)+(-2)=-4\ne-2.\] Indeed sufficiently negative indices obstruct this condition for any nonzero finite-dimensional filtered vector space. On a pure weight-\(n\) structure, the intersection formula defines Hodge summands on \(p+q=n\); it must not impose the stated off-diagonal condition on all filtration intersections. An equivalent correct purity criterion is \[V_\mathbb C=F^pV_\mathbb C\oplus\overline{F^{\,n-p+1}V_\mathbb C}
 \quad\hbox{for every }p .\]\par
\subsection*{Dependency and repair boundary}
The two local errors are assessed Minor. Example 3.31 does not require mixed-category semisimplicity. The correct homological dimension-one statement, different-weight Hom vanishing and higher-weight Ext vanishing suffice for its application. For the last assertion, an extension of pure \(X_i\) by pure \(X_j\) with \(j>i\) splits: strictness of \(W\) identifies \(W_iE\to X_i\) as an isomorphism, providing a section. Together with the correct pure-Hodge criterion, these properties repair the stated local reasoning without discarding mixed extension data. The entire realization theorem is not independently reconstructed here, and these examples do not demonstrate failure of a central motivic or \(p\)-adic conclusion.\par
\section*{Other checks and remaining limits}
The two triangles in Proposition 4.64 produce the chain in Remark 4.68 by composing their connecting maps. With \(A=\iota^!M\), \(B=\iota^*M\), \(C=\chi_{\rm alg}j^*M\), \(D=\Psi\operatorname{An}^*j^*M\) and \(E=D(-1)\), the two triangles are \(C\to D\to E\) and \(A\to B\to C\). Each adjacent composition in \[\cdots\longrightarrow B_n\longrightarrow D_n\longrightarrow E_n
 \longrightarrow A_{n-2}\longrightarrow B_{n-2}\longrightarrow\cdots\] vanishes by one of the original long exact sequences. Remark 4.69 gives \(A=M_k(X_k)\{-d-1\}\); setting \(n=-r\) turns \(A_{n-2}\) into \(H_{2d-r}(X_k)(-d-1)\). This verifies the two-step degree change and twist in Corollary 4.70. The dimension-zero good-reduction case gives the expected segment \(\mathbb Q\xrightarrow{\rm id}\mathbb Q\xrightarrow0\mathbb Q(-1)
\xrightarrow{\rm id}\mathbb Q(-1)\).\par
Corollary 4.70 has discrete valuation, flatness, properness, regularity, generic smoothness and polystable reduction among its hypotheses. Remark 4.71 conditions exactness on weight–monodromy. Construction of a chain complex was not misread as unconditional exactness. The more general introductory \(\mathcal O_C\) formulation was read, but its full descent bridge from the displayed body hypotheses was not independently established. This is retained as a scope limitation, without a manufactured counterexample. The generalization suggested in Example 4.73 remains explicitly an expectation.\par
No independent reconstruction of every cited foundational theorem, six-functor formalism, higher coherence, descent argument or weight–monodromy result is claimed. Reading all text and mechanically extracting all pages are recorded separately from the bounded proofs actually checked.\par
\section*{Conclusion and state}
The highest supported assessment is Major, for Problem Problem 1 (\#prob:purity)’s false premise in the central comparison derivation. Problems Problem 2 (\#prob:product) and Problem 3 (\#prob:hodge) have local impact on the evidence checked. The final comparison is not disproved; explicit local corrections have been described with their verification boundaries.\par
Item | Assessment | Repair checked\par
Tate purity premise | Major | Spectral condition and orthogonality\par
Product coalgebra | Minor | Actual comonad in the example\par
Hodge exposition | Minor | Extension and filtration criteria\par
Scientific adjudication | Pending | No independent acceptance\par
Author adoption | Unknown | No verified correction obtained\par
Human reviews | None recorded | No human attestation\par
9 F. Binda, M. Gallauer and A. Vezzani, Motivic monodromy and \(p\)-adic cohomology theories, Warwick WRAP accepted manuscript, record 188439, 51 PDF pages. DOI: `10.4171/jems/1634`. https://wrap.warwick.ac.uk/id/eprint/188439/. B. Drew, Rectification of Deligne’s mixed Hodge structures, arXiv:1511.08288v1, 26 November 2015, p. 9. This is an independently consulted reference, distinct from [Dre17] in the reviewed manuscript. https://arxiv.org/pdf/1511.08288. J. Francis and D. Gaitsgory, Chiral Koszul duality, author-hosted 43-page PDF, §3.5, p. 21; retrieved 27 October 2026. https://sites.math.northwestern.edu/\textasciitilde{}jnkf/writ/chiral.pdf. D. Gaitsgory and N. Rozenblyum, A study in derived algebraic geometry, Volume II, author-hosted PDF, Chapter 6, §2.7, printed p. 240 / PDF p. 273; retrieved 27 October 2026. This annotated file is not asserted byte-identical to the 2017 printing. https://people.mpim-bonn.mpg.de/gaitsgde/Book/Vol2.pdf.\par
\section*{Supporting checks}
The exact input PDF, original manifest, full extracted text, per-claim evidence, source-page images, coverage and findings ledgers are retained with this report. Rendering-time false inspection states remain preserved; later viewing is a separate fact.\par
Their actual scopes, limitations and a corrected reference attribution are retained. They are not human reviewers or independent acceptance executors. Model/UI observations supplied by the external dispatcher and the runtime’s self-description are separately recorded in provenance; no independent backend-identity attestation is made.\par
\end{document}
